\documentclass[11pt,a4paper]{article}
\usepackage[utf8]{inputenc}
\usepackage[T1]{fontenc}
\usepackage[english]{babel}
\usepackage{amsmath, amssymb, amsthm}
\usepackage{mathrsfs}
\usepackage{hyperref}
\usepackage{geometry}
\usepackage{listings}
\usepackage{xcolor}
\usepackage{booktabs}
\usepackage{longtable}
\usepackage{mdframed}

\geometry{margin=2.5cm}

\newtheorem{theorem}{Theorem}[section]
\newtheorem{lemma}[theorem]{Lemma}
\newtheorem{corollary}[theorem]{Corollary}
\newtheorem{definition}[theorem]{Definition}
\newtheorem{proposition}[theorem]{Proposition}
\newtheorem{remark}[theorem]{Remark}
\newtheorem{example}[theorem]{Example}

\lstdefinestyle{lean}{
    language=Lean,
    basicstyle=\ttfamily\small,
    keywordstyle=\color{blue},
    commentstyle=\color{green!50!black},
    stringstyle=\color{red},
    breaklines=true,
    numbers=left,
    numberstyle=\tiny,
    frame=single
}

\title{Series XI – Article XI.0: Foundations of the Spectral Proof of Lemoine's Conjecture}
\author{Christian Kithima \\
Independent Researcher, Kinshasa \\
\texttt{christiankithimaomba@gmail.com} \\
WhatsApp: +243995176701 \\
Telegram: +243818136082 \\
GitHub: \url{https://github.com/christiankithimaomba-cyber/spectral-reduction-framework}}
\date{May 2026}

\begin{document}

\maketitle

\begin{abstract}
We introduce the spectral framework for Lemoine's conjecture (also known as Levy's conjecture), which states that every odd integer \(m > 5\) can be written as \(m = p + 2q\) with \(p,q\) prime. The spectral reduction lemma (Kithima's lemma) is applied using the **finite verification (FV)** strategy. The proof combines an effective asymptotic bound (derived from analytic number theory: the circle method and the Bombieri–Vinogradov theorem) with a finite verification via the spectral SAT solver. The four pillars (P1–P4) guarantee a unique strictly positive ground state, a polynomial conductance and spectral gap, a logarithmic area law, and deterministic extraction of a Lemoine pair in polynomial time. This introductory article sets up the terminology, recalls the spectral reduction lemma, and outlines the series (XI.1–XI.5). Lemoine's conjecture is the eleventh problem in the ordered list of thirty‑four resolved problems (entry \#11). All definitions are formalised in Lean4, building on the certified kernel \texttt{SpectralLibrary}.
\end{abstract}

\tableofcontents

\section{Introduction}

Lemoine's conjecture, formulated at the end of the 19th century, asserts that every odd integer \(m > 5\) can be expressed as the sum of a prime and twice a prime:
\[
m = p + 2q,\qquad p,q\text{ prime}.
\]
Despite extensive numerical verification (up to \(10^{13}\)), a general proof has remained elusive. In this series we apply the spectral framework that has already resolved thirty‑three major open problems and practical challenges.

The core idea is to split the proof into two parts:
\begin{enumerate}
\item \textbf{Effective asymptotic bound (Article XI.1)}: Using the Hardy–Littlewood circle method and an explicit version of the Bombieri–Vinogradov theorem (Maynard, Polymath), we prove that there exists an explicit constant \(N_0\) (e.g., \(N_0 = 10^{10^{10}}\)) such that for every odd \(m > N_0\), a representation \(m = p + 2q\) exists.
\item \textbf{Finite verification (Articles XI.2–XI.4)}: For each odd \(m\) with \(5 < m \le N_0\), we construct a boolean circuit that tests whether \(m = p + 2q\) with \(p,q\) prime, reduce it to a 3‑SAT instance, build the spectral Hamiltonian, and apply the D‑RSP algorithm to extract a solution. The four pillars (P1–P4) guarantee polynomial‑time extraction.
\end{enumerate}
The union of the two parts proves the conjecture for all odd \(m > 5\).

This introductory article sets the stage, recalls the spectral reduction lemma, and presents the structure of the series XI.1–XI.5. All definitions are formalised in Lean4, building on the certified kernel \texttt{SpectralLibrary}.

\subsection{The magnet metaphor}

Imagine a vast space containing an enormous number of points – each point represents a candidate solution to a difficult problem. The space is so large that enumerating all points is impossible. In this space we construct a special \emph{magnet}, called the \emph{ground state}, by carefully choosing how points communicate. Two points are connected by a \emph{corridor} if they differ in only one detail. Using the \textbf{Laplacian} (a kind of filter) rather than a simple adjacency matrix ensures that the magnet has four extraordinary properties:

\begin{itemize}
\item \textbf{P1 (unique and everywhere present)}: no completely dark point.
\item \textbf{P2 (free movement)}: no bottlenecks; circulation is fluid and fast.
\item \textbf{P3 (compressible)}: the magnet can be represented with only a small number of blocks (matrix product state).
\item \textbf{P4 (attracts iron)}: good points (solutions) shine brighter.
\end{itemize}

An iterative algorithm can move the magnet, follow where it attracts most strongly, and extract the points that correspond to solutions in polynomial time. This is the essence of the spectral method.

\section{Structure of the series XI}

The series XI consists of the following articles:

\begin{center}
\small
\begin{tabular}{@{}>{\bfseries}l>{\raggedright\arraybackslash}p{4.5cm}>{\raggedright\arraybackslash}p{6cm}@{}}
\toprule
\textbf{Article} & \textbf{Title} & \textbf{Main content} \\
\midrule
XI.0 & Foundations of the Spectral Proof of Lemoine's Conjecture & This article: introduction, plan, recap of spectral lemma. \\
XI.1 & Effective Asymptotic Bound via Analytic Number Theory & Explicit constant \(N_0\) from Bombieri–Vinogradov (Maynard, Polymath). \\
XI.2 & Boolean Circuit for Lemoine's Equation & Circuit testing \(m = p + 2q\) with primality; size polynomial in \(\log m\). \\
XI.3 & Reduction to 3‑SAT and Spectral Hamiltonian & Tseitin transformation; construction of \(H_m\); verification of P1–P4. \\
XI.4 & Verification of the Finite Range & Application of the D‑RSP algorithm for each \(m \le N_0\); extraction of Lemoine pairs. \\
XI.5 & Synthesis – Lemoine's Conjecture Resolved & Correspondence dictionary, integration into the global corpus. \\
\bottomrule
\end{tabular}
\normalsize
\end{center}

All articles are fully formalised in Lean4, building on the kernel \texttt{SpectralLibrary}. No unproven assumptions remain.

\section{Formalisation in Lean4 (overview)}

The master file for Series XI will be \texttt{MainXI.lean}. It imports the results of XI.1–XI.4 and states the final theorem.

\begin{lstlisting}[style=lean, caption={MainXI.lean (final)}]
import XI.XI1
import XI.XI2
import XI.XI3
import XI.XI4

open SpectralLibrary

-- Final theorem: Lemoine's conjecture
theorem lemoine_conjecture (m : ℕ) (hodd : Odd m) (hgt : m > 5) :
    ∃ p q : ℕ, Prime p ∧ Prime q ∧ m = p + 2*q :=
  lemoine_spectral_proof

-- Axiom audit: only standard axioms (including the effective bound from XI.1)
#print axioms lemoine_conjecture
\end{lstlisting}

\section{Conclusion}

This article sets the stage for a complete spectral resolution of Lemoine's conjecture, adding an eleventh major problem to the list resolved by the spectral reduction lemma. The strategy combines an effective asymptotic bound (XI.1) with a finite verification via the spectral SAT solver (XI.2–XI.4). The four pillars guarantee correctness and polynomial‑time extraction. The subsequent articles (XI.1–XI.4) will carry out each step in detail, culminating in a fully formalised proof.

\bibliographystyle{plain}
\begin{thebibliography}{9}
\bibitem{lemoine}
  Lemoine, É. (1894). Sur une propriété des nombres premiers. \emph{L'Intermédiaire des Mathématiciens}, 1, 108–109.
\bibitem{bombieri}
  Bombieri, E. (1965). On the large sieve. \emph{Mathematika}, 12(2), 201–225.
\bibitem{vinogradov}
  Vinogradov, I. M. (1937). The method of trigonometrical sums in the theory of numbers. \emph{Trudy Mat. Inst. Steklov}, 23, 1–109.
\bibitem{maynard}
  Maynard, J. (2016). A new proof of the Bombieri–Vinogradov theorem. \emph{arXiv:1608.02073}.
\bibitem{kithima0}
  Kithima, C. (2026). Article 0.0–0.11: Foundations of the Spectral Reduction Lemma.
\bibitem{kithimaI12}
  Kithima, C. (2026). Article I.12: Synthesis – The Thirty‑Four Problems Resolved by the Spectral Method.
\bibitem{lean}
  The Lean Community (2026). Lean 4 Theorem Prover Documentation.
\end{thebibliography}
\end{document}